% !TeX root = ../main.tex
\chapter{Mô hình baseline}
\label{chap:baseline}

\section{Quy trình huấn luyện và kiểm thử}

Dataset được chia một lần theo tỷ lệ 80/20 với
\texttt{random\_state=42} và \texttt{stratify=y}. Tập training gồm 455 mẫu
(170 malignant, 285 benign); tập test gồm 114 mẫu (42 malignant, 72 benign).
Baseline là \texttt{DecisionTreeClassifier} với seed
\texttt{random\_state=42}. Các tham số cấu trúc giữ giá trị mặc định:
\texttt{criterion=gini}, \texttt{splitter=best} và
\texttt{max\_depth=None}. Mô hình không chứa tuning ngầm.

Tất cả mô hình cải thiện dùng chính split này. Việc chọn hyperparameter thực
hiện bằng 5-fold cross-validation trên 455 mẫu training. Đây là biện pháp tránh
data leakage do nhìn kết quả test nhiều lần trong quá trình tuning.

\section{Kết quả baseline}

\begin{table}[H]
\centering
\caption{Hiệu năng mô hình baseline trên tập test}
\label{tab:baseline-metrics}
\begin{tabular}{lc}
\toprule
Metric & Giá trị \\
\midrule
Training accuracy & 1,0000 \\
Testing accuracy & 0,9123 \\
Error rate & 0,0877 \\
Weighted precision & 0,9161 \\
Weighted recall & 0,9123 \\
Weighted F1-score & 0,9130 \\
Malignant precision & 0,8478 \\
Malignant recall & 0,9286 \\
Malignant F1-score & 0,8864 \\
\bottomrule
\end{tabular}
\end{table}

Accuracy test bằng 91,23\%, tương ứng error rate 8,77\%. Baseline phân loại
đúng 104/114 mẫu. Confusion matrix ở Hình~\ref{fig:baseline-cm} cho thấy mô
hình nhận đúng 39/42 ca malignant và 65/72 ca benign; có 3 ca malignant bị dự
đoán thành benign và 7 ca benign bị dự đoán thành malignant.
Các metric malignant trong báo cáo được tính riêng cho nhãn 0 theo quy ước của
dataset; vì vậy malignant recall đo trực tiếp tỷ lệ ca ác tính được phát hiện.

\begin{figure}[H]
\centering
\includegraphics[width=0.62\textwidth]{baseline_confusion_matrix.png}
\caption{Confusion matrix của Decision Tree baseline}
\label{fig:baseline-cm}
\end{figure}

\section{Cây kết quả}

Cây đầy đủ có depth 7, 19 leaf và 37 node. Vì hình đầy đủ khó đọc khi thu nhỏ,
Hình~\ref{fig:baseline-readable} trình bày bốn tầng đầu để phục vụ phân tích;
file \texttt{baseline\_tree\_full.png} được giữ trong gói kết quả kỹ thuật.

\begin{figure}[H]
\centering
\includegraphics[width=\textwidth]{baseline_tree_readable.png}
\caption{Cấu trúc dễ đọc của baseline (bốn tầng đầu, từ mức 0 đến mức 3)}
\label{fig:baseline-readable}
\end{figure}
